\documentclass{article}
\usepackage[slovene]{babel}


\usepackage{ragged2e}
\usepackage{lmodern}
\usepackage{graphicx}

\usepackage{subcaption}

\usepackage{enumitem}
% \setlist[itemize,1]{noitemsep, topsep=0pt, label=\heartsuit}
% \setlist[itemize,2]{noitemsep, topsep=0pt, label=\diamondsuit}
% \setlist[enumerate,1]{label=\Alph*.}
% \setlist[enumerate,2]{label*=\arabic*.}
% \setlist[description]{style=nextline}

\usepackage{array,booktabs}

\usepackage{multirow}


\usepackage{lipsum}

\begin{document}

% \tableofcontents

\setcounter{page}{0}

\lipsum[1-3]
\newpage

% \lipsum[1-3]
% \begin{Center}
%   \includegraphics[width=5cm]{./macka.jpeg}
% \end{Center}
% \lipsum[4-6] 
\lipsum[1-3]
\begin{figure}
  \centering
  \includegraphics[width=5cm]{./macka.jpeg}
  \caption{To je lepa mehiška mačka.}
\end{figure}
\lipsum[4-6] 

\clearpage


\vspace{1cm} \hrule \vspace{1cm}

\hrule
\section{To je section}
\subsection{To je subsection}
\subsubsection{To je subsubsection}
\paragraph{To je paragraph}
\subparagraph{To je subparagraph}


% \vspace{1cm} \hrule \vspace{1cm}
\section{Drugi razdelek}
\lipsum[4-6]
\subsection{a je to podrazdelek?}
\lipsum[7]
\subsection{Ja, to je podrazdelek!}
\lipsum[8-10]

\section{Privi razdelek}
\lipsum[1-3]
\subsection*{Podrazdelek brez številke}
\lipsum[11-13]

\section[Tretji razdelek]{Tretji razdelek z dolgim naslovom, ki se razteza čez več vrstic}
\lipsum[14-15]

\vspace{2cm}


\begin{FlushLeft}
  Ta besedilo je 
  levo poravnan.
  \LaTeX{} se ne trudi narediti
  vrstic z enakimi dolžinami.

  \lipsum[1][1-2]
\end{FlushLeft}

\begin{FlushRight}
  Ta besedilo je desno poravnano.
  Kot prej tudi tu vrstice nimajo
  enakih dolžin.

  \lipsum[1][1-2]
\end{FlushRight}

\begin{Center}
  V središču 
  sveta.

  \lipsum[1][1-2]
\end{Center}

\vspace{2cm}

\noindent 

\begin{enumerate}
  \item ena
  \item dva
  \item tri
\end{enumerate}

\begin{itemize}
  \item prva točka neurejena seznama.
  \item druga točka neurejena seznama.
\end{itemize}

\begin{description}
  \item[Prva:] točka opisa.
  \item[Druga:] točka opisa.
  \item[Tretja, z daljšim imenom:] točka opisa.
\end{description}

\vspace{2cm}

\begin{enumerate}
  \item Različna okolja lahko mešamo
  po lastnem okusu:
    \begin{itemize}
      \item Toda to lahko postane smešno.
      \item[-] To se začne s pomišljajem.
    \end{itemize}
  \item Zapomnite si:
    \begin{description}
      \item[Neumne] zadeve ne bodo postale pametne, 
      če so v seznamu.
      \item[Pametne] zadeve, za razliko, pa lahko čudovito 
      prikažemo s seznamom.
    \end{description}
\end{enumerate}


\clearpage

% Z uporabo možnosti \texttt{label=\textbackslash alph*)} lahko uporabljamo črke namesto številk:
\begin{enumerate}[label=\alph*)]
  \item Z možnostjo \texttt{label=\textbackslash textit\{(\textbackslash roman*)\}}, dobimo \textit{(i)}, \textit{(ii)}, \textit{(iii)} \dots
  \item Podobno, z možnostjo \texttt{label=\textbackslash textbf\{\textbackslash Alph*.\}}, dobimo \textbf{A.}, \textbf{B.}, \textbf{C.}...
\end{enumerate}

\vspace{2cm}
% Lahko spremenimo oznake neurejenih seznamov.
% \begin{itemize}[label=\heartsuit]
%   \item Tukaj uporabljamo \texttt{label=\textbackslash heartsuit}.
%   \item Z možnostjo \texttt{label=\textbackslash textendash} dobimo pomišljaj.
%   \item Z možnostjo \texttt{label=\textbackslash textasteriskCentered} dobimo zvezdico.
% \end{itemize}

\vspace{2cm}
% Z uporabo možnosti \texttt{style=nextline} lahko dosežemo, da se opisi začnejo v novi vrstici:
\begin{description}[style=nextline]
  \item[Jabolko] Rdeče ali zeleno sadje.
  \item[Pomaranča] Okroglo oranžno sadje.
\end{description}


\vspace{2cm}

\begin{enumerate}[label=\arabic*.]
  \item  Prva točka na prvi stopnji.
    \begin{enumerate}[label*=\Alph*.]
      \item Prva točka na drugi stopnji.
      \item Druga točka na drugi stopnji.
    \end{enumerate}
\end{enumerate}

\begin{enumerate}[label=\Roman*., resume]
  \item To je drugačen seznam, z rimskimi številkami.
  \item Številke pa se nadaljujejo iz prejšnjega seznama.
  \item Nikoli ne smete nadaljevati seznama z različnimi oznakami.
\end{enumerate}


\clearpage

\begin{itemize}
  \item Neurejeni seznami prve stopnje imajo kot oznake srčke.
  \item Opazujte, kako je bil razmik nastavljen tudi globalno.
  \begin{itemize}
    \item Neurejeni seznami na drugi stopnji imajo diamanti.
  \end{itemize}
\end{itemize}

\vspace{2cm}

\begin{enumerate}
  \item  Prav tako smo nastavili oznake urejenih seznamov.
  \begin{enumerate}
    \item Zdaj tisti na drugi stopnji privzeto podedujejo oznake (e.g. A.1.)
  \end{enumerate}
\end{enumerate}

\vspace{2cm}

\begin{description}
  \item[Opisni seznami] so nastavljeni tako, da se opisi začnejo v novi vrstici.
  \item[To je] zelo uporabno, kadar so opisi daljši.
\end{description}


\vspace{2cm}

Tabele v \LaTeX u se lahko poravnajo na sredino:
\begin{tabular}{ll} 1&2\\3&4 \end{tabular}
na dno: 
\begin{tabular}[b]{ll} 1&2\\3&4 \end{tabular}
ali na vrh: 
\begin{tabular}[t]{ll} 1&2\\3&4 \end{tabular}
glede na besedilo okoli njih.

\vspace{2cm}
\begin{tabular}{l|c r|}
  levo& središče & desno \\ \hline
  1 & 2 & 3 \\ \cline{1-1}
  4 & 5 & 6 \\ \cline{2-3}
  7 & 8 & 9 \\ \hline
\end{tabular}


\vspace{2cm}

\begin{tabular}{ l p{2cm} m{2cm} b{2cm}}
  \hline
  \texttt{l} & \texttt{p\{2cm\}} & \texttt{m\{2cm\}} & \texttt{b\{2cm\}} \\
  \hline
  Celica &
  Vrh besedila je poravnan s celico. &
  Središče besedila je poravnan s celico. &
  Dno besedila je poravnan s celico. \\\hline
\end{tabular}

\vspace{2cm}

\begin{tabular}{c p{.7\textwidth}}
  \toprule
  Znak & Opis \\
  \midrule
  \texttt{l} & Poravnava besedila v stolpcu na levo. \\
  \texttt{c} & Poravnava besedila v stolpcu na sredino. \\
  \texttt{r} & Poravnava besedila v stolpcu na desno. \\
  \texttt{p{<širina>}} & Besedilo v stolpcu bo zavito, če preseže določeno širino. \\
  \texttt{m{<širina>}} & Kot \texttt{p{<širina>}}, vendar bo besedilo navpično poravnano na sredino celice. \\
  \texttt{b{<širina>}} & Kot \texttt{p{<širina>}}, vendar bo besedilo navpično poravnano na dno celice. \\
  \bottomrule
  \end{tabular}

  \vspace{2cm}

  \begin{tabular}
    { l @{\hspace{1cm}}  r @{,} l @{ €}}
    % 
    \toprule
    Izdelek & \multicolumn{2}{c}{Cena}\\
    \midrule
    Espresso & 1&50 \\
    Kava z mlekom & 2&50 \\
    Čaj & 1&80 \\
    \bottomrule
  \end{tabular}
  

  \vspace{2cm}

  \begin{tabular}{ l c l }
  \toprule
  & \multicolumn{2}{c}{Reakcija} \\\cmidrule{2-3}
  Oseba & Obraz & Razpoloženje \\
  \midrule
  Jaz & :) & Vesel \\
  Vi & :| & Zaskrbljen \\
  Svoj partner & \multicolumn{2}{c}{Ni na voljo} \\
  \bottomrule
  \end{tabular}
\vspace{2cm}

  \begin{tabular}{ l c l }
  \toprule
 \multirow{2}{*}{Oseba} & \multicolumn{2}{c}{Reakcija} \\\cmidrule{2-3}
  & Obraz & Razpoloženje \\
  \midrule
  \multirow{2}{*}{Mi} & :) & Vesel \\
   & :| & Zaskrbljen \\
  Svoj partner & \multicolumn{2}{c}{Ni na voljo} \\
  \bottomrule
  \end{tabular}
\vspace{2cm}

  \begin{tabular}{ >{ \bfseries } l  >{\mdseries } p{0.6\textwidth} }
  \toprule
  Krepko besedilo & Navadno besedilo \\
  \midrule
  Krepko besedilo & Besedilo v prvem stolpcu je krepko. \\
  Krepko besedilo & Besedilo v prvem stolpcu je vedno krepko. \\
  Krepko besedilo & Res! Besedilo v prvem stolpcu je vedno krepko. \\
  \bottomrule
  \end{tabular}


\vspace{2cm}



\clearpage 
\includegraphics[width=2cm]{./macka.jpeg}
\includegraphics[height=2cm]{./macka.jpeg}
\includegraphics[width=2cm,height=2cm]{./macka.jpeg}

\vspace{2cm}

\includegraphics[scale=0.2]{./macka.jpeg}
\includegraphics[scale=0.1]{./macka.jpeg}
\includegraphics[scale=-0.1]{./macka.jpeg}

\clearpage

\lipsum[1][3] \includegraphics[width=1cm, angle=45]{./macka.jpeg} \lipsum[1][4]

% \vspace{.5cm}
% \hrule
% \vspace{.5cm}

\lipsum[1][3] \includegraphics[width=1cm, angle=180]{./macka.jpeg} \lipsum[1][4]

% \vspace{.5cm}
% \hrule
% \vspace{.5cm}

% \vspace{1cm}
\lipsum[1][3] \includegraphics[width=1cm, angle=180, origin=c]{./macka.jpeg} \lipsum[1][4]

% \vspace{.5cm}
% \hrule
% \vspace{.5cm}

% \vspace{1cm}
\lipsum[1][3] \includegraphics[width=1cm, angle=180, origin=lt]{./macka.jpeg} \lipsum[1][4]


\clearpage 


\lipsum[1-3]
\begin{Center}
  \includegraphics[width=5cm]{./macka.jpeg}
\end{Center}
\lipsum[4-6]  

\vspace{2cm}

\begin{tabular}{ l c r }
  \texttt{l} & \texttt{c} & \texttt{r} \\
  levo  &
  središče &
  desno  \\
  1 & 2 & 3 \\
\end{tabular}


\clearpage

\begin{figure}
  \centering
  \begin{subfigure}[t]{0.4\textwidth}
    \includegraphics[width=\textwidth]{./macka.jpeg}
    \caption{Zelo dooooooooooooolga pojaznila za najlepšo Mačko na svetu.}
  \end{subfigure}
  \hfill
  \begin{subfigure}[t]{0.4\textwidth}
    \includegraphics[width=\textwidth]{./pes.jpg}
    \caption{Kuža}
  \end{subfigure}
  \caption{Pojaznilo za obe podsliki}
\end{figure}


\end{document}